\section{Domain-Namen}

\subsection{Grundlagen und Aufbau}

\begin{definition}
    \textbf{Domain-Name:} Weltweit eindeutige, menschenlesbare Adresse zur Identifikation und Adressierung von Rechnern und Diensten im Internet (wird über das DNS in IP-Adressen aufgelöst).
\end{definition}

\begin{remark}
    \textbf{Hierarchischer Aufbau:}
    \begin{itemize}
        \item \textbf{Top-Level-Domain (TLD):} Oberste Ebene der Namenshierarchie:
              \begin{itemize}
                  \item \textit{ccTLD (country code):} Länderspezifisch (z.B. \texttt{.ch}, \texttt{.de})
                  \item \textit{gTLD (generic):} Generisch (z.B. \texttt{.com}, \texttt{.org}, \texttt{.swiss})
              \end{itemize}
        \item \textbf{Second-Level-Domain (SLD):} Der eigentliche gewählte Name (z.B. \texttt{ethz} in \texttt{ethz.ch}).
        \item \textbf{Subdomains:} Weitere Untergliederungen (z.B. \texttt{moodle.ethz.ch}).
    \end{itemize}
\end{remark}

\subsection{Vergabesystem und Akteure}

\begin{remark}
    \textbf{Akteure bei der Domain-Vergabe:}
    \begin{itemize}
        \item \textbf{ICANN:} Globale Non-Profit-Organisation, koordiniert das Domain Name System und die Root-Zone weltweit.
        \item \textbf{Nationale Ordnung (Registry):} Verwaltet die Datenbank der jeweiligen TLD (z.B. SWITCH bzw. BAKOM für \texttt{.ch}).
        \item \textbf{Händler (Registrar):} Akkreditierte Verkäufer (z.B. Hostpoint, Infomaniak), über die Endkunden Domains registrieren.
        \item \textbf{Registrant:} Der Endkunde bzw. Domain-Halter.
    \end{itemize}
\end{remark}

\subsection{Rechtliche Natur und Vergabe}

\begin{definition}
    \textbf{First come, first served:} Domain-Namen werden nach dem Prioritätsprinzip vergeben: Wer eine freie Domain zuerst anmeldet, erhält den Zuschlag. Die Registrierungsstellen führen keine Vorprüfung auf allfällige Marken- oder Namensrechte durch.
\end{definition}

\begin{important}
    \textbf{Dienstleistung, nicht Eigentum!} \\
    An einem Domain-Namen erwirbt man \textbf{kein Sacheigentum} (kein Eigentum im Sinne von Art. 641 ZGB). \\
    Es handelt sich um eine rein \textbf{vertragliche Dienstleistung}: Ein schuldrechtliches, relatives Nutzungsrecht auf Eintragung und Routing im DNS gegen Bezahlung einer periodischen Gebühr.
\end{important}